\documentclass[14pt]{extarticle}
\usepackage{fontspec}
\usepackage{polyglossia}
\usepackage{amsmath}
\usepackage{amsfonts}
\usepackage{amssymb}
\usepackage{pgfplots}
\pgfplotsset{compat=1.18}

% Hebrew setup
\setdefaultlanguage{hebrew}
\setotherlanguage{english}
\newfontfamily\hebrewfont[Script=Hebrew]{Arial}

% Math font setup for proper Greek symbols
\usepackage{unicode-math}
\setmathfont{Latin Modern Math}

% Increase line spacing throughout document
\linespread{1.3}

\usepackage[margin=1in]{geometry}

% Left-aligned equation environment
\newenvironment{lefteqs}{%
    \begin{english}%
    \begin{flushleft}%
    \renewcommand{\arraystretch}{1.3}%
    $\begin{array}{@{}l@{}}%
}{%
    \end{array}$%
    \end{flushleft}%
    \end{english}%
}

\begin{document}

\begin{center}
    {\LARGE \textbf{ממ"ן 11 - פיזיקה קוונטית 1}}\\[0.5em]
    {\large שי רואימי - 318986270}\\[0.3em]
\end{center}

\hrule
\vspace{2em}

\textbf{\large שאלה 1}

\textbf{א)}
התבקשנו למצוא את צפיפות ההסתברות $\rho$ ואת צפיפות הזרם $J$, עבור:

\begin{english}
    \begin{flalign*}
        (1) \quad & \Psi(x,t)=A \cdot e^{-(kx-\omega t)^2}       & \\
        (2) \quad & \Psi(x,t)=A \cdot e^{-a^2x^2+i(kx-\omega t)} &
    \end{flalign*}
\end{english}

כאשר
$J(x,t)=\frac{i \hbar}{2m}(\Psi \frac{\partial \Psi^*}{\partial x} - \Psi^* \frac{\partial \Psi}{\partial x})$
,ו $\rho(x,t)=|\Psi(x,t)|^2$.

עבור (1):

\begin{lefteqs}
    \rho(x,t) = A^2 e^{-2(kx-\omega t)^2} \\
    \dfrac{\partial \Psi}{\partial x} = \dfrac{\partial \Psi^*}{\partial x} \\
    \Psi=\Psi^* \\
    J(x,t) = 0
\end{lefteqs}

עבור (2):

\begin{lefteqs}
    \Psi^* = Ae^{-a^2x^2 -i(kx-\omega t)} \\
    \rho(x,t) = \Psi \Psi^* = Ae^{-a^2x^2+i(kx-\omega t)} \cdot Ae^{-a^2x^2-i(kx-\omega t)}
    \rho(x,t)=A^2 e^{-2a^2x^2} \\
    \dfrac{\partial \Psi}{\partial x} = (-2a^2x+ik)Ae^{-a^2x^2+i(kx-\omega t)} \\
    \dfrac{\partial \Psi^*}{\partial x} = (-2a^2x-ik)Ae^{-a^2x^2 -i(kx-\omega t)} \\
    J(x,t)=\frac{i\hbar}{2m} \cdot ((-2a^2x-ik)|\Psi|^2 - (-2a^2x+ik)|\Psi|^2) = \frac{\hbar k}{m} |\Psi|^2 \\
    J(x,t) = \frac{\hbar k}{m} A^2e^{-2a^2x^2}
\end{lefteqs}

\textbf{ב)}
עבור (1)

\begin{lefteqs}
    \frac{\partial \rho}{\partial t} = A^2 \cdot 4\omega(kx-wt)e^{-2k(kx-\omega t)^2}
\end{lefteqs}

ואילו
\begin{lefteqs}
    \frac{\partial J}{\partial x} = 0
\end{lefteqs}
משוואת הרציפות אינה מתקיימת

עבור (2)
\begin{lefteqs}
    \rho(x,t)=A^2e^{-2a^2x^2} \Rightarrow \dfrac{\partial \rho}{\partial t} = 0 \\
    \dfrac{\partial J}{\partial x} = \frac{\hbar k}{m} A^2 \cdot (-4a^2x)e^{-2a^2x^2} = -\frac{4\hbar k a^2 x}{m} A^2 e^{-2a^2x^2}
\end{lefteqs}

משוואת הרציפות מתקיימת רק עבור $x=0$.


\vspace{3em}
\hrule
\vspace{2em}

\textbf{\large שאלה 2}

\textbf{א)}
חלקיק נע במימד אחד וברגע $t_0$ מתקיים $\Psi(x,t_0)=N\cdot e^{\frac{-x^2}{2a^2}}$

נמצא את מקדם הנרמול $N$. מקדם הנרמול מקיים: $\int_{-\infty}^{\infty} |\Psi(x,t_0)|^2 dx = 1$.

\begin{lefteqs}
    \int_{-\infty}^{\infty}N e^{\frac{-x^2}{2a^2}} \cdot N e^{\frac{-x^2}{2a^2}} dx = 1 \\
    N^2 \int_{-\infty}^{\infty} e^{\frac{-x^2}{a^2}} dx = 1 \\
    N^2 \cdot 2 \int_{0}^{\infty} e^{\frac{-x^2}{a^2}} dx = 1 \\
    N^2 \cdot 2 \cdot \frac{a \sqrt{\pi}}{2} = 1 \\
    N^2 \cdot a \sqrt{\pi} = 1 \\
    N^2 = \frac{1}{a \sqrt{\pi}} \\
    N = \sqrt[4]{\frac{1}{a^2 \pi}}
\end{lefteqs}

\textbf{ב)}
אי הוודאות במקום הוא $\sigma_x$, כאשר:

\begin{lefteqs}
    \sigma_x = \sqrt{\langle x^2 \rangle - \langle x \rangle^2} \\
    \langle x \rangle = \int_{-\infty}^{\infty} x \cdot N^2 e^{\frac{-x^2}{a^2}} dx = 0 \\
\end{lefteqs}

האינטגרל שווה ל0 כי הפונקציה היא זוגית והכפלה בx הופכת אותה לאי זוגית.

אינטגרל של פונקציה אי זוגית מעל טווח סימטרי סביב 0 שווה ל0.

\begin{lefteqs}
    \langle x^2 \rangle = \int_{-\infty}^{\infty} x^2 \cdot N^2 e^{\frac{-x^2}{a^2}} dx = 2 \int_{0}^{\infty} x^2 \cdot N^2 e^{\frac{-x^2}{a^2}} dx \\
    =2N^2 \cdot \int_{0}^{\infty} x^2 e^{\frac{-x^2}{a^2}} dx = 2N^2 \cdot \sqrt{\pi} \cdot 2 \cdot (\frac{a}{2})^3 \\
    \langle x^2 \rangle = \frac{a^2}{2} \\
    \sigma_x = \sqrt{\frac{a^2}{2} - 0} = \frac{a}{\sqrt{2}}
\end{lefteqs}

\textbf{ג)}
\begin{lefteqs}
    \sigma_p = \sqrt{\langle p^2 \rangle - \langle p \rangle^2} \\
    \langle p \rangle = \int_{-infty}^{\infty} \Psi^* \cdot (-i\hbar \frac{\partial}{\partial x}) \cdot \Psi dx = \int_{-\infty}^{\infty} N e^{\frac{-x^2}{2a^2}} \cdot (-i\hbar \frac{\partial}{\partial x}) \cdot N e^{\frac{-x^2}{2a^2}} dx \\
    = \int_{-\infty}^{\infty} N e^{\frac{-x^2}{2a^2}} \cdot (-i\hbar) \cdot N \cdot (-\frac{x}{a^2}) e^{\frac{-x^2}{2a^2}} dx \\
    = \int_{-\infty}^{\infty} N^2 \cdot \frac{i\hbar x}{a^2} e^{\frac{-x^2}{a^2}} dx = 0
\end{lefteqs}
גם פה האינטגרל שווה 0 בתור אינטגרל של פונקציה אי זוגית בטווח סימטרי סביב 0
\begin{lefteqs}
    \langle p^2 \rangle = \int_{-\infty}^{\infty} \Psi^* \cdot (-i\hbar \frac{\partial}{\partial x})^2 \cdot \Psi dx = \int_{-\infty}^{\infty} N e^{\frac{-x^2}{2a^2}} \cdot (-i\hbar \frac{\partial}{\partial x})^2 \cdot N e^{\frac{-x^2}{2a^2}} dx \\
    \dfrac{\partial \Psi}{\partial x} = -\frac{x}{a^2} \cdot N e^{\frac{-x^2}{2a^2}}
    \dfrac{\partial^2 \Psi}{\partial x^2} = (-\frac{1}{a^2} + \frac{x^2}{a^4}) \cdot N e^{\frac{-x^2}{2a^2}} \\
    \langle p^2 \rangle = \int_{-\infty}^{\infty} N^2 \hbar^2 \cdot \frac{1}{a^2} e^{\frac{-x^2}{a^2}} dx - \int_{-\infty}^{\infty} N^2 \hbar^2 \cdot \frac{x^2}{a^4} e^{\frac{-x^2}{a^2}} dx \\
    = \frac{N^2 \hbar^2}{a^2} ( \int_{-\infty}^{\infty} e^{\frac{-x^2}{a^2}} dx - \frac{1}{a^2}\int_{-\infty}^{\infty} x^2 e^{\frac{-x^2}{a^2}} dx ) \\
    = \frac{2 \cdot N^2 \hbar^2}{a^2} (\frac{a \sqrt{\pi}}{2} - \frac{1}{a^2} \cdot \frac{a^3 \sqrt{\pi}}{4}) \\
    \langle p^2 \rangle = \frac{\hbar^2}{2a^2} \\
    \sigma_p = \sqrt{\frac{\hbar^2}{2a^2} - 0} = \frac{\hbar}{\sqrt{2}a} \\
    \sigma_x \sigma_p = \frac{a}{\sqrt{2}} \cdot \frac{\hbar}{\sqrt{2}a} = \frac{\hbar}{2} \geq \frac{\hbar}{2} \\
\end{lefteqs}
עקרון אי הוודאות מתקיים.

\vspace{3em}
\hrule
\vspace{2em}

\textbf{\large שאלה 3}

\textbf{א)}
\begin{lefteqs}
    \Psi(x,0) =
    \begin{cases}
        A\left(\frac{x}{a}\right),         & 0 \le x \le a    \\[6pt]
        A\left(\frac{b - x}{b - a}\right), & a \le x \le b    \\[6pt]
        0,                                 & \text{otherwise}
    \end{cases}
\end{lefteqs}
דרישת הנרמול היא  $\int_{-\infty}^{\infty} |\Psi(x,0)|^2 dx = 1$.
\begin{lefteqs}
    \int_{-\infty}^{\infty} \Psi \Psi* dx = \int_{a}{b} (A\frac{b-x}{b-a})^2 dx + \int_{0}^{a} (A\frac{x}{a})^2 dx \\
    = \frac{A^2}{(b-a)^2}\int_{a}^{b} b^2-2bx+x^2 dx + \frac{A^2}{a^2} \int_{0}^{a} x^2 dx \\
    = \frac{A^2}{(b-a)^2}\left[ b^2x-bx^2+\frac{x^3}{3} \right]_{a}^{b} + \frac{A^2}{a^2} \left[\frac{x^3}{3} \right]_{0}^{a} \\
    = \frac{A^2}{(b-a)^2} (b^3 - b^3 + \frac{b^3}{3} - ab^2 + a^2b - \frac{a^3}{3}) +  \frac{A^2 a^3}{3a^2}  \\
    = \frac{A^2(b-a)^3}{3(b-a)^2} + \frac{A^2 a}{3} = \frac{A^2 (b-a)}{3} + \frac{A^2 a}{3} = \frac{A^2 (b-a+a)}{3} = \frac{A^2 b}{3} \\
    \frac{A^2 b}{3} = 1 \Rightarrow A = \sqrt{\frac{3}{b}}
\end{lefteqs}

\textbf{ב)}

\begin{tikzpicture}
    \begin{axis}[
            axis lines=middle,
            xmin=0, xmax=6,
            ymin=0, ymax=3,
            xlabel={$x$},
            ylabel={$\Psi(x)$},
            xtick={2,5},          % positions of a and b
            xticklabels={$a$,$b$},
            ytick={2},            % peak height
            yticklabels={$\sqrt{\frac{3}{b}}$},
        ]

        % Triangle
        \addplot[thick] coordinates {
                (0,0)
                (2,2)   % (a, A)
                (5,0)   % (b, 0)
            };

    \end{axis}
\end{tikzpicture}

\textbf{ג)}
הפונקציה $\Psi(x,t)$ ממשית אי  שלילית בתחום $[0,b]$ ומקבלת מקסימום בנקודה $x=a$ לכן גם $|\Psi(x,t)|^2$ מקבלת מקסימום בנקודה $x=a$. לכן, הסיכוי למצוא את החלקיק בנקודה $x=a$ הוא הגבוה ביותר.

\textbf{ד)}
\begin{lefteqs}
    \int_{0}^{a} |\Psi(x,0)|^2 dx = \int_{0}^{a} \frac{A^2}{a^2} x^2 dx = \frac{A^2}{a^2} \left[ \frac{x^3}{3} \right]_{0}^{a}
    = \frac{A^2}{a^2}\frac{a^3}{3} = \frac{3}{b} \frac{a}{3} = \frac{a}{b} \\
    b = a \Rightarrow \int_{0}^{a} |\Psi(x,0)|^2 dx = 1 \\
    b = 2a \Rightarrow \int_{0}^{a} |\Psi(x,0)|^2 dx = \frac{1}{2}
\end{lefteqs}

\textbf{ה)}
\begin{lefteqs}
    \langle x \rangle = \int_{-\infty}^{\infty} x |\Psi(x,0)|^2 dx = \int_{a}{b} x(A\frac{b-x}{b-a})^2 dx + \int_{0}^{a} x(A\frac{x}{a})^2 dx \\
    = \frac{A^2}{(b-a)^2}  \int_{a}^{b}  x(b^2-2bx+x^2) dx + \frac{A^2}{a^2} \int_{0}^{a} x^3 dx \\
    = \frac{A^2}{(b-a)^2} \left[ \frac{b^2x^2}{2} - \frac{2bx^3}{3} + \frac{x^4}{4} \right]_{a}^{b} + \frac{A^2}{a^2} \left[ \frac{x^4}{4} \right]_{0}^{a} \\
    = \frac{A^2}{(b-a)^2} \left( \frac{b^4}{2} - \frac{2b^4}{3} + \frac{b^4}{4} - \frac{a^2 b^2}{2} + \frac{2ba^3}{3} - \frac{a^4}{4} \right) + \frac{A^2 a^2}{4} \\
    = \frac{A^2}{(b-a)^2} \left( \frac{b^4}{12} - \frac{a^2 b^2}{2} + \frac{2ba^3}{3} - \frac{a^4}{4} \right) + \frac{A^2 a^2}{4} \\
    = \frac{12}{12(b-a)^2} (b^4  -3a^4  +8ba^3 -6b^2a^2) + \frac{A^2a^2}{4} \\
    = \frac{1}{12} \frac{3}{b} \frac{1}{(b-a)^2} (b^4  -3a^4  +8ba^3 -6b^2a^2) + \frac{3}{4b} a^2 \\
    = \frac{a}{4b(b-a)^2}(b-a)^2(b^2+2ab-3a^2) + \frac{3a^2}{4b} \\
    = \frac{b^2+2ab-3a^2}{4b}+\frac{3a^2}{4b} = \frac{b^2+2ab}{4b}  =\frac{b+2a}{4} \\
    \langle x \rangle =  \frac{b+2a}{4}
\end{lefteqs}


\vspace{3em}
\hrule
\vspace{2em}

\textbf{\large שאלה 4}

לפי משוואת  שרודינגר:
\begin{lefteqs}
    \frac{\partial \Psi}{\partial t} = \frac{i\hbar}{2m} \frac{\partial^2 \Psi}{\partial x^2} -\frac{i}{\hbar} V \Psi \\
    \frac{\partial \Psi^*}{\partial t} = -\frac{i\hbar}{2m} \frac{\partial^2 \Psi^*}{\partial x^2} + \frac{i}{\hbar} V \Psi^* \\
    \Rightarrow \frac{d}{dt} \int_{-\infty}^{\infty} \Psi_{1}^* \Psi_{2} dx = \int_{-\infty}^{\infty} \frac{\partial}{\partial t} (\Psi_{1}^* \Psi_{2}) dx  =  \int_{-\infty}^{\infty} (\frac{\partial  \Psi_{1}^*}{\partial t}\Psi_{2} + \Psi_{1}^* \frac{\partial \Psi_{2}}{\partial t}) dx \\
    = \int_{-\infty}^{\infty} \left( -\frac{i\hbar}{2m}\cdot \frac{\partial^2 \Psi_{1}^*}{\partial x^2} + \frac{i}{\hbar}V\Psi_{1}^* \right)\Psi_{2} + \left( \frac{i\hbar}{2m} \cdot  \frac{\partial^2 \Psi_{2}}{\partial x^2} - \frac{i}{\hbar} V \Psi_{2} \right) \Psi_{1}^* dx \\
    = \frac{i\hbar}{2m} \int_{-\infty}^{\infty} \frac{\partial^2 \Psi_{2}}{\partial x^2} \Psi_{1}^* - \frac{\partial^2 \Psi_{1}^*}{\partial x^2} \Psi_{2} dx \\
    = \frac{i\hbar}{2m} \left( \left[ \Psi_{1}^*\frac{\partial \Psi_2}{\partial x} \right]_{-\infty}^{\infty}  
    - \int_{-\infty}^{\infty} \frac{\partial \Psi_{1}^*}{\partial x} \cdot \frac{\partial \Psi_2}{\partial x} dxi  - \left[ \frac{\partial \Psi_{1}^*}{\partial x} \Psi_2 \right]_{-\infty}^{\infty} + \int_{-\infty}^{\infty} \frac{\partial \Psi_{1}^*}{\partial x} \cdot \frac{\partial \Psi_2}{\partial x} dx \right) = 0   
\end{lefteqs}

\vspace{3em}
\hrule
\vspace{2em}

\textbf{\large שאלה 5}

\textbf{א)}
\begin{lefteqs}
    \Psi(x,0) = \begin{cases}
        A(a^2 - x^2), & -a \le x \le a    \\[6pt]
        0,            & \text{otherwise}
    \end{cases} \\
    \int_{-\infty}^{\infty} | \Psi(x,0)|^2 dx = \int_{-a}^{a} A^2 (a^2 - x^2)^2 dx = A^2 \int_{-a}^{a} (a^4 - 2a^2x^2 + x^4) dx \\
    = A^2 \left[ a^4x - \frac{2a^2x^3}{3} + \frac{x^5}{5} \right]_{-a}^{a} = A^2 \cdot \frac{16 a^5}{15} \\
    A = \frac{1}{4}\sqrt{\frac{15}{a^5}}
\end{lefteqs}

\textbf{ב)}
\begin{lefteqs}
    \langle x \rangle = \int_{-\infty}^{\infty} x |\Psi(x,0)|^2 dx = \int_{-a}^{a} x A^2 (a^2 - x^2)^2 dx = 0
\end{lefteqs}
אינטגרל של פונקציה אי זוגית בטווח סימטרי סביב 0 שווה ל0.

\textbf{ג)}
\begin{lefteqs}
    \langle p \rangle = \int_{-\infty}^{\infty} \Psi^* \cdot (-i\hbar \frac{\partial}{\partial x}) \cdot \Psi dx
    = i\hbar \int_{-a}^{a} A^2 (a^2 - x^2) \cdot (-2x) dx \\
    = -2A^2 i\hbar \int_{-a}^{a} x(a^2 - x^2) dx = 0
\end{lefteqs}

\textbf{ד)}
\begin{lefteqs}
    \langle x^2 \rangle = \int_{-\infty}^{\infty} x^2 |\Psi(x,0)|^2 dx = \int_{-a}^{a} x^2 A^2 (a^2 - x^2)^2 dx  \\
    = 2A^2 \int_{0}^{a} x^2(x^4-2a^2x^2+a^4) dx =2A^2 \int_{0}^{a} (x^6 - 2a^2x^4 + a^4x^2) dx \\
    = 2A^2 \left[ \frac{x^7}{7} - \frac{2a^2x^5}{5} + \frac{a^4x^3}{3} \right]_{0}^{a} = 2A^2 \cdot \frac{8 a^7}{105} = \frac{a^2}{7} \\
    \Rightarrow \langle x^2 \rangle = \frac{a^2}{7}
\end{lefteqs}

\textbf{ה)}
\begin{lefteqs}
    \langle p^2 \rangle = \int_{-\infty}^{\infty} \Psi^* \cdot (-i\hbar \frac{\partial}{\partial x})^2 \cdot \Psi dx = -\hbar^2 \int_{-a}^{a} A(a^2 - x^2) \cdot (-2A) dx \\
    = 2A^2\hbar^2 \cdot 2 \int_{0}^ {a}a^2-x^2 dx  = 4\hbar^2A^2 \left[ a^2x-\frac{x^3}{3} \right]_{0}^{a} \\
    = 4\hbar^2A^2 (a^3-\frac{a^3}{3})=  4\hbar^2  \cdot \frac{15}{16a^5}\cdot  \frac{2a^3}{3}  = \frac{5\hbar^2}{2a^2} \\
    \Rightarrow \langle p^2 \rangle = \frac{5\hbar^2}{2a^2}
\end{lefteqs}

\textbf{ו)}
\begin{lefteqs}
    \sigma_x = \sqrt{\langle x^2 \rangle - \langle x \rangle^2} = \sqrt{\frac{a^2}{7}} = \frac{a}{\sqrt{7}} \\
\end{lefteqs}

\textbf{ז)}
\begin{lefteqs}
    \sigma_p = \sqrt{\langle p^2 \rangle - \langle p \rangle^2} = \sqrt{\frac{5\hbar^2}{2a^2}} = \frac{\hbar}{a} \sqrt{\frac{5}{2}} \\
\end{lefteqs}

\textbf{ח)}
\begin{lefteqs}
    \sigma_x \sigma_p = \frac{a}{\sqrt{7}} \cdot \frac{\hbar}{a} \sqrt{\frac{5}{2}} = \hbar \sqrt{\frac{5}{14}} \ge \frac{\hbar}{2}
\end{lefteqs}
עקרון אי הוודאות מתקיים.

\end{document}